\documentclass{article} % For LaTeX2e
\usepackage{iclr2026_conference,times}

\input{math_commands.tex}

\usepackage{hyperref}
\usepackage{url}
\usepackage{enumitem}

\iclrfinalcopy

\title{Membership Inference Attack on a ResNet-18 Image Classifier}

\author{Sahil Malhan \\
\texttt{sahil.malhan00@gmail.com}}

\begin{document}

\maketitle
\vspace{-1.1em}

\noindent\textbf{GitHub Repository:}
\url{https://github.com/mayankmalhan123/-Membership-Inference-Attack_Assignment_1}

\section{Introduction}

This assignment studies membership inference against a trained image classifier. The attacker receives a pretrained ResNet-18, a public dataset with known membership labels, and a private dataset whose membership values are hidden. The goal is to assign each private sample a score in $[0,1]$ that reflects how likely it is that the sample was part of the target model's training set. Membership inference attacks are a standard way to probe privacy leakage from machine learning models \citep{shokri2017mia,carlini2022firstprinciples}. In this task, the evaluation metric is TPR@5\%FPR, so the attack must rank likely members above non-members while keeping false positives low.

\section{Approach}

I treated the problem as a supervised attack-learning task. Rather than retraining the target classifier, I used \texttt{pub.pt} as labeled data for training a separate attack model. The key idea is that members and non-members may look slightly different through the outputs of the target model, even when they are sampled from the same underlying distribution.

I first tested simple threshold baselines based on true-label confidence and negative cross-entropy loss. These are common first-step attacks because training samples often have lower loss and higher confidence \citep{carlini2022firstprinciples,zarifzadeh2023lowcost}. However, on the public data these one-dimensional scores were close to random for the assignment metric, which suggested that the leakage signal was weak and needed richer features.

I therefore built a feature vector from the target model outputs for each sample. The feature vector contains raw logits, softmax probabilities, true-label confidence, per-sample loss, prediction entropy, correctness, and a one-hot encoding of the ground-truth class. These features capture confidence, calibration, uncertainty, and class-specific behavior. On top of them, I trained a small multilayer perceptron with one hidden layer and binary cross-entropy loss. I kept the attack model intentionally simple in order to reduce overfitting and make the behavior of the method easier to interpret.

\section{Validation And Results}

I evaluated the attack by splitting \texttt{pub.pt} into attack-train and attack-validation subsets while preserving both class balance and the member/non-member ratio within each class. I then measured the same leaderboard-oriented metric used by the assignment.

The strongest simple baselines, such as true-label confidence and negative loss, achieved only about 0.0507 TPR at 5\% FPR on the public data, which is effectively near-random for this metric. In contrast, the learned attack improved the validation score to about 0.0756 TPR at 5\% FPR, with a ROC-AUC of about 0.5246. This shows that the target model leaks some membership information, but that the signal is subtle and is better captured by combining several weak signals rather than relying on a single threshold.

\noindent\textbf{Best public leaderboard score:} \texttt{TO\_BE\_FILLED\_AFTER\_FIRST\_SUBMISSION}

\section{Conclusion}

This assignment demonstrates that privacy leakage can persist even when the attacker only has access to a trained model and a set of candidate samples. A successful membership inference attack means that participation in the training set can itself become sensitive information. In practice, this matters whenever the presence of a record in training reveals something private about a user, such as medical status, service usage, or membership in a protected group. Even though the attack here is not extremely strong, it still shows that standard predictive models can leak nontrivial information about their training data, which is why privacy auditing and privacy-aware training remain important in trustworthy machine learning.

\bibliography{references}
\bibliographystyle{iclr2026_conference}

\end{document}
